\section{Пораждащи функции}

Пораждащите функции представляват полезна трансформация на случайните величини. Подобно на архитектурна скица, която дава различен поглед върху една сграда в зависимост от ъгъла, пораждащата функция предоставя мощен аналитичен инструмент за изследване на свойствата на дискретните разпределения.

\begin{keydefn}
Нека $X : \Omega \rightarrow \mathbb{N} \cup \{0\}$ е целочислена случайна величина (т.е. възможните ѝ стойности са $0, 1, 2, \dots$). Тогава под \textbf{пораждаща функция} (probability generating function) на $X$ разбираме функцията $g_X(s)$, дефинирана като:
\[ g_X(s) := \E s^X = \sum_{n=0}^{\infty} s^n \mathbb{P}(X=n), \quad \text{за } s \in [-1, 1]. \]
\end{keydefn}

Този ред кодира цялото разпределение на $X$: коефициентът пред $s^n$ е точно $\mathbb{P}(X=n)$. Записът с безкраен ред е удобен и когато $X$ приема само краен брой стойности, защото тогава всички останали коефициенти просто са нула. Така на едно разпределение се съпоставя една аналитична функция, от която впоследствие могат да се извлекат отново отделните вероятности. Именно това е смисълът на аналогията със скиците на една сграда: различното представяне прави видими различни свойства на същия обект.

С помощта на пораждащата функция могат лесно да се пресмятат математическото очакване и дисперсията на случайната величина, както и да се възстановят нейните вероятности.

\begin{prop}
Нека $X$ е целочислена случайна величина с пораждаща функция $g_X(s)$. Тогава:
\begin{enumerate}
    \item[а)] $\E X = g_X'(1)$
    \item[б)] $\Var X = g_X''(1) + g_X'(1) - (g_X'(1))^2$
    \item[в)] $k! \mathbb{P}(X=k) = g_X^{(k)}(0)$, за $k \ge 0$.
\end{enumerate}
\end{prop}


\begin{proof}
а) По дефиниция $g_X(s) = \sum_{n=0}^{\infty} s^n \mathbb{P}(X=n)$. Диференцирайки по $s$, получаваме:
\[ g_X'(s) = \frac{d}{ds} \sum_{n=0}^{\infty} s^n \mathbb{P}(X=n) = \sum_{n=1}^{\infty} n s^{n-1} \mathbb{P}(X=n) = \E[X s^{X-1}]. \]
Замествайки $s=1$, намираме $g_X'(1) = \sum_{n=1}^{\infty} n \mathbb{P}(X=n) = \E X$.

б) Аналогично, втората производна е:
\[ g_X''(s) = \E[X(X-1)s^{X-2}]. \]
При $s=1$ получаваме $g_X''(1) = \E[X(X-1)] = \E X^2 - \E X$.
От формулата за дисперсия $\Var X = \E X^2 - (\E X)^2$, можем да изразим $\E X^2 = g_X''(1) + \E X = g_X''(1) + g_X'(1)$. Така:
\[ \Var X = g_X''(1) + g_X'(1) - (g_X'(1))^2. \]

в) За да намерим вероятността $X$ да приеме стойност $k$, диференцираме $g_X(s)$ $k$ пъти. Всички членове със степен по-малка от $k$ се анулират, а членът със степен $k$ става $k! \mathbb{P}(X=k)$. При заместване с $s=0$, всички останали членове (съдържащи $s$) се нулират, и остава точно $k! \mathbb{P}(X=k)$.
\end{proof}

\begin{cor}
Нека $X$ и $Y$ са целочислени случайни величини. Тогава $X \stackrel{d}{=} Y$ (равни по разпределение) тогава и само тогава, когато $g_X \equiv g_Y$.
\end{cor}

\begin{proof}
Ако $X\stackrel{d}{=}Y$, то $\mathbb{P}(X=n)=\mathbb{P}(Y=n)$ за всяко $n\ge0$, следователно съответните коефициенти в двата реда съвпадат и $g_X=g_Y$. Обратно, ако $g_X=g_Y$, то и производните им във всяка степен съвпадат. От вече доказаната формула
\[
k!\,\mathbb{P}(X=k)=g_X^{(k)}(0)=g_Y^{(k)}(0)=k!\,\mathbb{P}(Y=k)
\]
следва равенство на всички точкови вероятности, а значи и равенство по разпределение.
\end{proof}


Важно свойство на пораждащите функции се проявява при суми от независими случайни величини. 

\begin{defn}
Дискретните случайни величини $X_1, \dots, X_n$ (или $(X_j)_{j=1}^\infty$) имат \textbf{еднакво разпределение}, ако $X_i \stackrel{d}{=} X_j$ за всеки $i, j$.
\end{defn}


\begin{prop}
Нека $X_1, X_2, \dots, X_n$ са независими в съвкупност целочислени случайни величини. Ако $Y = \sum_{j=1}^n X_j$, то:
\[ g_Y(s) = \prod_{j=1}^n g_{X_j}(s). \]
Ако в допълнение $X_j$ са еднакво разпределени, то $g_Y(s) = (g_{X_1}(s))^n$.
\end{prop}

\begin{proof}
Ще го покажем за $n=2$. Нека $Y = X_1 + X_2$. Тогава:
\[ g_Y(s) = \E s^{X_1+X_2} = \sum_{j=0}^{\infty} \sum_{i=0}^{\infty} s^{i+j} \mathbb{P}(X_1=j, X_2=i). \]
Поради независимостта на $X_1$ и $X_2$, съвместната вероятност се разпада: $\mathbb{P}(X_1=j, X_2=i) = \mathbb{P}(X_1=j)\mathbb{P}(X_2=i)$.
\[ g_Y(s) = \sum_{j=0}^{\infty} s^j \mathbb{P}(X_1=j) \sum_{i=0}^{\infty} s^i \mathbb{P}(X_2=i) = g_{X_1}(s) g_{X_2}(s). \]
\end{proof}

Тук независимостта е съществена: тя позволява съвместната вероятност да се разложи на произведение. Ползата е, че прякото намиране на разпределението на една сума изисква да се обходят всички комбинации от стойности на събираемите, които дават един и същ резултат. При пораждащите функции тази комбинаторна задача се заменя с умножение на функции. Ако събираемите са и еднакво разпределени, произведението се свежда до една степен.

\begin{supp}[Пораждащата функция и функцията на разпределение]\label{supp:pgf-cdf}
Освен вероятностите и моментите, пораждащата функция кодира и функцията на разпределение:
за целочислена $X \ge 0$, при приетата тук строга конвенция $F_X(n) = \mathbb{P}(X < n)$,
\[ \sum_{n \ge 0} s^n F_X(n) = \frac{s\, g_X(s)}{1-s}, \qquad |s| < 1 . \]
Делението на $1-s$ съответства на натрупването: умножаването на реда по
$\frac{1}{1-s} = 1 + s + s^2 + \dots$ заменя всяка вероятност със сумата на всички преди
нея, което е точно преходът от $\mathbb{P}(X=n)$ към $F_X(n)$. Допълнителният множител
$s$ отчита строгото неравенство: той измества реда с една позиция, така че $\mathbb{P}(X=n)$
да не участва във $F_X(n)$. При широката конвенция $F_X(n) = \mathbb{P}(X \le n)$ множителят
отпада.
\end{supp}

\section{Основни дискретни разпределения}

Пораждащите функции ни дават единен апарат, с който да изведем очакването, дисперсията и вероятностите на изброените по-долу разпределения. Всички те се разглеждат върху схемата на Бернули.

\subsection{Схема на Бернули}

Схемата на Бернули е базов вероятностен модел, описващ поредица от независими експерименти с два възможни изхода (успех или неуспех) и постоянна вероятностна структура.
Двете изисквания трябва да се различават. Независимостта означава, че резултатът от един опит не променя вероятностите в останалите; еднаквото разпределение означава, че вероятността $p$ не се сменя от опит към опит. Например последователно хвърляне на различни монети може да даде независими резултати, но не и еднакво разпределени, ако монетите имат различни вероятности за ези.
Разглеждаме редица $(X_j)_{j=1}^{\infty}$, които са независими в съвкупност и еднакво разпределени:
\begin{align*}
\mathbb{P}(X_j=1) &= p \quad (\text{вероятност за успех}) \\
\mathbb{P}(X_j=0) &= 1-p = q \quad (\text{вероятност за неуспех})
\end{align*}

\medskip\noindent\textbf{Разпределение на Бернули.}
Случайната величина $X_1$ се нарича разпределена по Бернули, $X_1 \sim \Ber(p)$.
\[
\begin{array}{c|c|c}
X_1 & 0 & 1 \\ \hline
\mathbb{P} & q & p
\end{array}
\]
Пораждащата ѝ функция е:
\[ g_{X_1}(s) = q \cdot s^0 + p \cdot s^1 = q + ps. \]
Математическото очакване и дисперсията са:
\[ \E X_1 = p, \quad \Var X_1 = \E X_1^2 - (\E X_1)^2 = p - p^2 = pq. \]

\subsection{Биномно разпределение}

Биномното разпределение описва броя на успехите в първите $n$ експеримента от схема на Бернули. Записваме $X \sim \Bin(n, p)$.
Случайната величина $X$ може да се представи като сума от независими бернулиеви величини: $X = \sum_{j=1}^n X_j$.

\begin{prop}
Нека $X \sim \Bin(n, p)$. Тогава:
\begin{enumerate}
    \item[а)] $g_X(s) = (q+ps)^n$
    \item[б)] $\E X = np$
    \item[в)] $\Var X = npq$
    \item[г)] $\mathbb{P}(X=k) = \binom{n}{k} p^k q^{n-k}$, за $0 \le k \le n$.
\end{enumerate}
\end{prop}

\begin{lecfigure}[Биномното разпределение при $n=10$. Върхът стои около $\E X = np$:
при $p = 0{,}3$ това е $3$, при $p = 0{,}5$ — $5$. Симетрия има само при $p = \frac12$; за
$p < \frac12$ разпределението е изтеглено надясно, защото отляво го притиска границата
$k \ge 0$. Височините във всеки панел се сумират до
единица.\label{fig:binomial-pmf}]
\begin{tikzpicture}[x=0.52cm, y=9.5cm]
 \foreach \i/\lab/\dat in {%
   0/{$p = 0{,}3$}/{0/0.028,1/0.121,2/0.233,3/0.267,4/0.200,5/0.103,6/0.037,7/0.009,8/0.001,9/0.000,10/0.000},
   1/{$p = 0{,}5$}/{0/0.001,1/0.010,2/0.044,3/0.117,4/0.205,5/0.246,6/0.205,7/0.117,8/0.044,9/0.010,10/0.001}}
 {
  \begin{scope}[xshift=\i*7.2cm]
    \foreach \k/\h in \dat {
      \fill[defback] ({\k-0.36},0) rectangle ({\k+0.36},\h);
      \draw[defframe, line width=0.4pt] ({\k-0.36},0) rectangle ({\k+0.36},\h);}
    \draw[axisline] (-0.8,0) -- (10.9,0) node[below right, font=\small] {$k$};
    \foreach \k in {0,2,4,6,8,10} { \node[below, font=\scriptsize] at (\k,-0.004) {\k}; }
    \node[below, font=\small] at (5,-0.028) {\lab};
  \end{scope}}
\end{tikzpicture}
\end{lecfigure}

\begin{proof}
а) Тъй като $X = \sum_{j=1}^n X_j$ и $X_j$ са независими и еднакво разпределени, то $g_X(s) = (g_{X_1}(s))^n = (q+ps)^n$.

б) $\E X = g_X'(1) = np(q+ps)^{n-1} \big|_{s=1} = np(p+q)^{n-1} = np$.

в) Дисперсията се получава по същия начин чрез втората производна; сметката е оставена на читателя.

г) $\mathbb{P}(X=k)$ може да се намери чрез $k$-тата производна в нулата:
\[ k! \mathbb{P}(X=k) = g_X^{(k)}(0) = \left. \frac{d^k}{ds^k} (q+ps)^n \right|_{s=0} = n(n-1)\dots(n-k+1) p^k q^{n-k}. \]
Следователно $\mathbb{P}(X=k) = \frac{n(n-1)\dots(n-k+1)}{k!} p^k q^{n-k} = \binom{n}{k} p^k q^{n-k}$.
\end{proof}

Като приближен модел от лекцията беше разгледан потокът от хора, завръщащи се в страната в началото на пандемията. Ако за всеки пристигащ означим със стойност $1$ заразен човек и със стойност $0$ незаразен, приемем една и съща вероятност за заразяване $p$ и независимост между пристигащите, то всеки индикатор е бернулиев. Броят заразени сред първите $n$ пристигнали е сумата на тези индикатори и следователно има биномно разпределение. Лекторът изрично представи това като грубо приближение: еднаквото $p$ предполага сходен регионален риск, а независимостта изключва общ източник на заразяване.

\begin{supp}[Най-вероятна стойност (мода)]\label{supp:mode}
Естествен въпрос е къде разпределението достига максимума си. Отговорът се получава, като
сравним два съседни члена: разглеждаме отношението $\mathbb{P}(X=k)/\mathbb{P}(X=k-1)$ и
търсим къде то минава през $1$ — дотам вероятностите растат, след това намаляват.
За биномното разпределение
\[
\frac{\mathbb{P}(X=k)}{\mathbb{P}(X=k-1)} = \frac{n-k+1}{k}\cdot\frac{p}{q} \ \ge\ 1
\iff k \le p(n+1),
\]
откъдето модата е $k^* = \lfloor p(n+1) \rfloor$. По същия начин за Поасоновото
разпределение $k^* = \lfloor \lambda \rfloor$. Забележете, че модата и очакването са
близки, но по принцип различни числа.
\end{supp}

\subsection{Геометрично разпределение}

Геометричното разпределение моделира броя на неуспехите до настъпване на първия успех в схема на Бернули. Записваме $X \sim \Geo(p)$. Формално:
\[ X = \min\left\{n \ge 1 : \sum_{j=1}^n X_j = 1\right\} - 1. \]
В някои източници геометричното разпределение брои всички опити до първия успех и тогава не се изважда единица. В лекцията е избрана конвенцията за броя на неуспехите, така че възможните стойности започват от $0$. Например редицата $0,0,0,1$ дава $X=3$.

\begin{prop}
Нека $X \sim \Geo(p)$. Тогава:
\begin{enumerate}
    \item[а)] $\mathbb{P}(X=k) = q^k p$, за $k \ge 0$.
    \item[б)] $g_X(s) = \frac{p}{1-qs}$
    \item[в)] $\E X = \frac{q}{p}$, $\Var X = \frac{q}{p^2}$
\end{enumerate}
\end{prop}

\begin{proof}
а) Събитието $\{X=k\}$ означава, че първите $k$ експеримента са неуспешни, а $k+1$-вият е успешен. Поради независимостта:
\[ \mathbb{P}(X=k) = \mathbb{P}(X_1=0) \dots \mathbb{P}(X_k=0) \mathbb{P}(X_{k+1}=1) = q^k p. \]

б) За пораждащата функция имаме:
\[ g_X(s) = \sum_{n=0}^{\infty} s^n \mathbb{P}(X=n) = \sum_{n=0}^{\infty} s^n q^n p = p \sum_{n=0}^{\infty} (qs)^n = \frac{p}{1-qs}. \]

в) Чрез диференциране на $g_X(s)$ получаваме
\[
g_X'(s)=\frac{pq}{(1-qs)^2},
\qquad
g_X''(s)=\frac{2pq^2}{(1-qs)^3}.
\]
Следователно
\[
\E X=g_X'(1)=\frac{pq}{(1-q)^2}=\frac{q}{p}
\]
и, понеже $1-q=p$,
\begin{align*}
\Var X
  &=g_X''(1)+g_X'(1)-\bigl(g_X'(1)\bigr)^2 \\
  &=\frac{2q^2}{p^2}+\frac{q}{p}-\frac{q^2}{p^2}
    =\frac{q^2+pq}{p^2}
    =\frac{q}{p^2}.
\end{align*}
Така очакването и дисперсията се получават с пряка аналитична сметка от пораждащата функция, без да се пресмятат поотделно съответните степенни редове.
\end{proof}

Интерпретацията не зависи от това кой от двата изхода е наречен „успех“. При честна монета броят тури преди първото ези е геометричен с $p=1/2$ и има очакване $1$. В модела с вероятност за заразяване $p=0{,}05$ средният брой незаразени преди първия заразен е $0{,}95/0{,}05=19$. При баскетболист, който вкарва наказателен удар с вероятност $0{,}9$, броят вкарани кошове преди първия пропуск също е геометричен, но тук „успехът“, който прекратява броенето, е именно пропускът и неговата вероятност е $0{,}1$.

\medskip\noindent\textbf{Свойство безпаметност (Memoryless property).}
Геометричното разпределение е единственото дискретно разпределение, което „няма памет“. Тоест, фактът, че вече са настъпили определен брой неуспехи, не променя вероятността за бъдещите експерименти.

\begin{prop}[Безпаметност]
Ако $X \sim \Geo(p)$, то за $\forall m \ge 0$ и $\forall k \ge 0$:
\[ \mathbb{P}(X \ge m+k \given X \ge k) = \mathbb{P}(X \ge m). \]
\end{prop}

\begin{proof}
\[ \mathbb{P}(X \ge m+k \given X \ge k) = \frac{\mathbb{P}(X \ge m+k \text{ и } X \ge k)}{\mathbb{P}(X \ge k)} = \frac{\mathbb{P}(X \ge m+k)}{\mathbb{P}(X \ge k)}. \]
Тъй като $\mathbb{P}(X \ge k) = q^k$, получаваме $\frac{q^{m+k}}{q^k} = q^m = \mathbb{P}(X \ge m)$.
\end{proof}

Една интерпретация е животът на компонент: ако процесор вече е работил $k$ дни, при модел с безпаметност вероятността да работи още $m$ дни е същата като за току-що поставен процесор. Лекторът противопостави това на човешкия живот, при който възрастта очевидно променя оставащото време и геометричният модел не е подходящ.

Същото свойство обяснява защо историята на независими хазартни опити не помага за следващия залог. Ако три пъти поред се е паднало черно или тура, условната вероятност редицата да продължи още два пъти е същата, както без тази предходна информация. Записването на старите резултати не променя бъдещите вероятности; това важи и за миналите тиражи на „6 от 49“ при модела на независими тегления.

\subsection{Отрицателно биномно разпределение}

Отрицателното биномно разпределение обобщава геометричното. То брои броя на неуспехите до настъпването на $r$-тия успех в схема на Бернули. Означаваме го с $X \sim \NegBin(r, p)$, където $r \ge 1$ е цяло число.
Забележете, че $\NegBin(1, p) = \Geo(p)$.
Случайната величина може да се представи като сума от независими геометрични величини: $X = \sum_{j=1}^r Y_j$, където $Y_j \sim \Geo(p)$. В лекцията това е свързано с броя на пристигналите пътници преди да бъдат установени $r$ заразени: интересува ни не първият установен случай, а моментът, в който броят им достига капацитета, за който системата е подготвена.
Разлагането на сумата има конкретен смисъл: $Y_1$ брои неуспехите до първия успех, след него броенето започва отново и $Y_2$ брои неуспехите до втория успех, и така до $Y_r$. Поради независимостта на опитите тези блокове са независими и имат едно и също геометрично разпределение. Тази структура е причината очакването, дисперсията и пораждащата функция да се получават непосредствено от съответните величини за геометричното разпределение.

\begin{prop}
Нека $X \sim \NegBin(r, p)$. Тогава:
\begin{enumerate}
    \item[а)] $g_X(s) = \left( \frac{p}{1-qs} \right)^r$
    \item[б)] $\E X = r\frac{q}{p}$, $\Var X = r\frac{q}{p^2}$
    \item[в)] $\mathbb{P}(X=k) = \binom{r+k-1}{k} p^r q^k$, за $k \ge 0$.
\end{enumerate}
\end{prop}

\begin{proof}
От представянето като сума и независимостта на $Y_1,\dots,Y_r$ следва
\[
g_X(s)=\prod_{j=1}^r g_{Y_j}(s)
      =\left(\frac{p}{1-qs}\right)^r.
\]
Линейността на очакването и адитивността на дисперсията за независими величини дават
\[
\E X=\sum_{j=1}^r\E Y_j=r\frac{q}{p},
\qquad
\Var X=\sum_{j=1}^r\Var(Y_j)=r\frac{q}{p^2}.
\]

За вероятностите използваме $k$-тата производна на пораждащата функция:
\begin{align*}
k!\,\mathbb{P}(X=k)
  &=g_X^{(k)}(0) \\
  &=p^r q^k r(r+1)\dots(r+k-1).
\end{align*}
След деление на $k!$ получаваме
\[
\mathbb{P}(X=k)
  =\frac{r(r+1)\dots(r+k-1)}{k!}p^rq^k
  =\binom{r+k-1}{k}p^rq^k.
\]
Това е аналитичният извод от лекцията; пораждащата функция свежда комбинаторния коефициент до формално диференциране.
\end{proof}

\subsection{Поасоново разпределение}

Поасоновото разпределение често се използва за моделиране на събития, настъпващи рядко в непрекъснато време или пространство (напр. брой катастрофи, брой клиенти за единица време).
Като други начални модели бяха посочени броят голове в мач и броят видими звезди в участък от небето. Тук разпределението се задава направо чрез вероятностите си, а не като конкретен брой успехи в схема на Бернули. По-дълбокото му извеждане чрез процеси със стационарни и независими нараствания беше изрично оставено извън курса. Практическата връзка, която предстои, е друга: при подходящ режим Поасоновото разпределение приближава биномното и избягва тежките пресмятания с биномни коефициенти.

\begin{defn}
Случайната величина $X$ е разпределена по Поасон с параметър $\lambda > 0$, $X \sim \Poi(\lambda)$, ако вероятностите ѝ са дадени от:
\[ \mathbb{P}(X=k) = \frac{\lambda^k}{k!} e^{-\lambda}, \quad \text{за } k \ge 0. \]
\end{defn}

Коректността на това разпределение следва от развитието на експоненциалната функция в ред на Тейлър:
\[ \sum_{k=0}^{\infty} \mathbb{P}(X=k) = e^{-\lambda} \sum_{k=0}^{\infty} \frac{\lambda^k}{k!} = e^{-\lambda} e^{\lambda} = 1. \]

\begin{prop}
Нека $X \sim \Poi(\lambda)$. Тогава $g_X(s) = e^{\lambda s - \lambda}$, и $\E X = \Var X = \lambda$.
\end{prop}

\begin{proof}
За пораждащата функция имаме:
\[ g_X(s) = \sum_{n=0}^{\infty} s^n \frac{\lambda^n}{n!} e^{-\lambda} = e^{-\lambda} \sum_{n=0}^{\infty} \frac{(s\lambda)^n}{n!} = e^{-\lambda} e^{\lambda s} = e^{\lambda s - \lambda}. \]
Диференцирайки, получаваме $g_X'(s) = \lambda e^{\lambda s - \lambda}$, което при $s=1$ дава $\E X = \lambda$.
Втората производна е $g_X''(s) = \lambda^2 e^{\lambda s - \lambda}$, откъдето $\Var X = \lambda^2 + \lambda - (\lambda)^2 = \lambda$.
\end{proof}

\section{Задачи}

\begin{enumerate}
    \item Доказателството на формулата
    \[
    g_{X_1+\dots+X_n}(s)=\prod_{j=1}^n g_{X_j}(s)
    \]
    беше направено за две независими случайни величини. Обобщете го за произволен краен брой събираеми.
    \item За $X\sim\Bin(n,p)$ довършете оставеното пресмятане с втората производна на $g_X(s)=(q+ps)^n$ и изведете $\Var X=npq$.
    \item За геометрично разпределена величина пресметнете $\E X$ направо от дефиниционния ред, без да използвате пораждащата функция.
    \item За $X\sim\Poi(\lambda)$ намерете $\E X$ и $\Var X$ чрез първите две производни на $g_X(s)=e^{\lambda(s-1)}$.
\end{enumerate}
